\section{Core Claim and Scope}

The central claim of \SPMPC{} is that open-liquid local planning should be treated as a predictive state problem, not only as a smooth-command generation problem.
For a mobile base carrying an open container, the future free-surface response depends on the current liquid displacement, liquid velocity, and residual oscillation phase.
Two command sequences with comparable velocity, acceleration, and smoothness statistics can therefore induce different liquid responses if they interact differently with the current liquid state.

\SPMPC{} addresses this issue at the online local-planning layer.
Given a geometrically feasible reference path, the planner repeatedly solves a finite-horizon \MPCC{} problem from the current robot, path-progress, and liquid modal state, then executes the first base command.
The method preserves the interface of an ordinary mobile-base local planner while making liquid dynamic memory part of the prediction model.

The technical body consists of three linked elements.
First, the robot state and path-progress state are augmented with a low-order slosh modal state.
Second, the \MPCC{} objective jointly penalizes path error, insufficient progress, speed-reference mismatch, control effort, control variation, and predicted slosh response.
Third, an optional Slosh-risk Reference Governor contracts the speed reference through short-horizon slosh prediction before the \MPCC{} solve.

The scope is intentionally bounded.
The method is not a high-fidelity free-surface simulator, a hard spill-prevention controller, or a closed-loop liquid-surface feedback controller.
Model-derived slosh quantities are used as prediction and diagnostic proxies.
Claims about real liquid motion must be supported by external observations such as RGB liquid-surface measurements or liquid-level sensing.
